\documentclass[12pt,a4paper]{article}
\usepackage[utf8]{inputenc}
\usepackage[T1]{fontenc}
\usepackage[vietnamese,english]{babel}
\usepackage{amsmath}
\usepackage{amsfonts}
\usepackage{amssymb}
\usepackage{amsthm}
\usepackage{algorithm}
\usepackage{algpseudocode}
\usepackage{graphicx}
\usepackage{hyperref}
\usepackage{xcolor}
\usepackage{listings}
\usepackage{tikz}
\usetikzlibrary{shapes,arrows,positioning,calc,decorations.pathreplacing}
\usepackage{booktabs}
\usepackage{geometry}
\usepackage{float}
\usepackage{longtable}
\geometry{margin=2.5cm}

\theoremstyle{definition}
\newtheorem{definition}{Definition}[section]
\newtheorem{theorem}{Theorem}[section]
\newtheorem{lemma}{Lemma}[section]
\newtheorem{example}{Example}[section]

\lstset{
    language=Python,
    basicstyle=\ttfamily\small,
    keywordstyle=\color{blue},
    commentstyle=\color{green!60!black},
    stringstyle=\color{red},
    numbers=left,
    numberstyle=\tiny,
    breaklines=true,
    frame=single,
    showstringspaces=false
}

\title{\textbf{Adapting Incremental MaxSAT from RTSS to CVRP}\\[1em]
\large A Detailed Technical Guide}
\author{Technical Documentation}
\date{\today}

\begin{document}

\maketitle

\tableofcontents
\newpage

%=============================================================================
\section{Introduction}
%=============================================================================

\subsection{Motivation}

The Real-Time Taxi-Sharing Service (RTSS) problem and the Capacitated Vehicle Routing Problem (CVRP) share fundamental similarities:
\begin{itemize}
    \item Both involve routing multiple vehicles
    \item Both have capacity constraints
    \item Both require optimization of total distance/cost
    \item Both are NP-hard combinatorial optimization problems
\end{itemize}

The RTSS paper by Zha et al. introduces an \textbf{Incremental MaxSAT} approach that efficiently solves the problem by:
\begin{enumerate}
    \item Encoding the problem using \textbf{connection variables} $l(i,j)$ and \textbf{reachability variables} $r(i,j)$
    \item Using \textbf{lazy constraint checking} for capacity/deadline constraints
    \item Applying \textbf{incremental SAT solving} with the RC2 solver
\end{enumerate}

This document details how to adapt this methodology to CVRP.

\subsection{Key Differences}

\begin{table}[H]
\centering
\begin{tabular}{|l|p{5cm}|p{5cm}|}
\hline
\textbf{Aspect} & \textbf{RTSS} & \textbf{CVRP} \\
\hline
Depot & Multiple (current taxi positions) & Single central depot \\
\hline
Route Structure & Open path (no return to start) & Closed circuit (return to depot) \\
\hline
Constraints & Capacity + Deadline & Capacity only \\
\hline
Customer Type & Pick-up and Drop-off pairs & Single service points \\
\hline
Ordering & Pick-up must precede drop-off & No pairing constraints \\
\hline
Time Dimension & Real-time, dynamic & Static, offline \\
\hline
\end{tabular}
\caption{Comparison of RTSS and CVRP}
\end{table}

%=============================================================================
\section{RTSS Methodology Overview}
%=============================================================================

\subsection{Boolean Variable Encoding}

The RTSS approach uses two types of Boolean variables:

\begin{definition}[Connection Variable $l(i,j)$]
$l(i,j) = \text{True}$ if and only if the vehicle travels \textbf{directly} from location $i$ to location $j$ (i.e., $(i,j)$ is an edge in the route).
\end{definition}

\begin{definition}[Reachability Variable $r(i,j)$]
$r(i,j) = \text{True}$ if and only if location $i$ is visited \textbf{before} location $j$ in the route.
\end{definition}

\textbf{Key Insight:} The relationship between $l$ and $r$ is:
\begin{equation}
l(i,j) \Rightarrow r(i,j)
\end{equation}

If we use edge $(i,j)$, then $i$ must be visited before $j$.

\subsection{Hamiltonian Path Constraints}

The RTSS paper encodes Hamiltonian path constraints using reachability variables with the following logical laws:

\subsubsection{Implication Rule}
\begin{equation}
l(i,j) \Rightarrow r(i,j)
\end{equation}

\subsubsection{Transitivity Law}
\begin{equation}
r(a,b) \wedge r(b,c) \Rightarrow r(a,c)
\end{equation}

\subsubsection{Chain Law}
\begin{equation}
r(a,b) \wedge r(b,c) \Rightarrow \neg l(c,a)
\end{equation}
(Prevents ``jumping back'' over an intermediate node)

\subsubsection{Acyclic Law}
\begin{equation}
\neg(r(a,b) \wedge r(b,a))
\end{equation}
(Ordering is asymmetric)

\subsection{Incremental Solving Strategy}

\begin{algorithm}[H]
\caption{RTSS Incremental Solving (solveRTSS)}
\begin{algorithmic}[1]
\Require WCNF formula with soft/hard clauses, RC2 solver
\Ensure Feasible solution or UNSAT

\State $model \gets$ rc2.compute()
\While{$model \neq$ None}
    \State $route \gets$ decodeModel($model$)
    \State $violation, nogood \gets$ checkExCondition($route$)
    
    \If{NOT $violation$}
        \State \Return $model$ \Comment{Found feasible solution}
    \Else
        \State rc2.add\_clause($nogood$) \Comment{Add nogood and continue}
        \State $model \gets$ rc2.compute()
    \EndIf
\EndWhile
\State \Return UNSAT
\end{algorithmic}
\end{algorithm}

The key function \texttt{checkExCondition} checks:
\begin{enumerate}
    \item \textbf{Capacity violations}: Total demand exceeds vehicle capacity
    \item \textbf{Deadline violations}: Arrival time exceeds deadline
\end{enumerate}

When a violation is found, a \textbf{nogood clause} is generated that says ``don't use this exact route configuration again''.

%=============================================================================
\section{Adapting to CVRP}
%=============================================================================

\subsection{Variable Mapping}

For CVRP with $n$ customers and $m$ vehicles:

\begin{table}[H]
\centering
\begin{tabular}{|c|c|p{7cm}|}
\hline
\textbf{RTSS} & \textbf{CVRP} & \textbf{Description} \\
\hline
$l_k(i,j)$ & conNet$[v][i][j]$ & Edge from $i$ to $j$ on vehicle $v$ \\
\hline
$r_k(i,j)$ & rchNet$[v][i][j]$ & $i$ before $j$ on vehicle $v$ \\
\hline
taxi $k$ & vehicle $v$ & Vehicle index $v \in \{0, \ldots, m-1\}$ \\
\hline
via points & customers & Customer index $i \in \{1, \ldots, n\}$ \\
\hline
pick/drop & single node & Each customer is a single service point \\
\hline
\end{tabular}
\caption{Variable Mapping from RTSS to CVRP}
\end{table}

\subsection{Modified Constraints}

\subsubsection{Depot Constraints (Different from RTSS)}

In CVRP, vehicles start and end at the \textbf{same depot} (node 0):

\begin{equation}
\sum_{j=1}^{n} \text{conNet}[v][0][j] = 1, \quad \forall v \in \{0, \ldots, m-1\}
\tag{C1: Leave depot once}
\end{equation}

\begin{equation}
\sum_{i=1}^{n} \text{conNet}[v][i][0] = 1, \quad \forall v \in \{0, \ldots, m-1\}
\tag{C2: Return to depot once}
\end{equation}

\subsubsection{Flow Conservation}

Each customer is visited exactly once by exactly one vehicle:

\begin{equation}
\sum_{v=0}^{m-1} \sum_{i=0, i \neq j}^{n} \text{conNet}[v][i][j] = 1, \quad \forall j \in \{1, \ldots, n\}
\tag{C3: One entry per customer}
\end{equation}

\begin{equation}
\sum_{v=0}^{m-1} \sum_{k=0, k \neq j}^{n} \text{conNet}[v][j][k] = 1, \quad \forall j \in \{1, \ldots, n\}
\tag{C4: One exit per customer}
\end{equation}

\subsubsection{Vehicle Consistency Constraint (NEW)}

If a customer is visited by vehicle $v$, both incoming and outgoing edges must belong to the SAME vehicle:

\begin{equation}
\text{conNet}[v][i][j] \wedge \text{conNet}[v'][j][k] \Rightarrow \text{False}, \quad \forall v \neq v'
\tag{C5: Vehicle consistency}
\end{equation}

In CNF: $\neg\text{conNet}[v][i][j] \vee \neg\text{conNet}[v'][j][k]$

This ensures that a customer's incoming and outgoing edges are on the same vehicle.

\subsubsection{Reachability Constraints (Same as RTSS)}

\textbf{Implication:}
\begin{equation}
\text{conNet}[v][i][j] \Rightarrow \text{rchNet}[v][i][j]
\end{equation}

In CNF: $\neg\text{conNet}[v][i][j] \vee \text{rchNet}[v][i][j]$

\textbf{Transitivity:}
\begin{equation}
\text{rchNet}[v][a][b] \wedge \text{rchNet}[v][b][c] \Rightarrow \text{rchNet}[v][a][c]
\end{equation}

In CNF: $\neg\text{rchNet}[v][a][b] \vee \neg\text{rchNet}[v][b][c] \vee \text{rchNet}[v][a][c]$

\textbf{Chain Law:}
\begin{equation}
\text{rchNet}[v][a][b] \wedge \text{rchNet}[v][b][c] \Rightarrow \neg\text{conNet}[v][c][a]
\end{equation}

In CNF: $\neg\text{rchNet}[v][a][b] \vee \neg\text{rchNet}[v][b][c] \vee \neg\text{conNet}[v][c][a]$

\textbf{No Back-and-Forth:}
\begin{equation}
\neg(\text{conNet}[v][i][j] \wedge \text{conNet}[v][j][i])
\end{equation}

In CNF: $\neg\text{conNet}[v][i][j] \vee \neg\text{conNet}[v][j][i]$

\subsection{Variable Reduction Trick}

\textbf{Key optimization from RTSS:}

For customers $i < j$ on the same vehicle, if $r(i,j) = \text{True}$, then $r(j,i) = \text{False}$.

Therefore, we only need ONE variable for each unordered pair:
\begin{align}
\text{rchNet}[v][i][j] &= \text{var\_id} \quad \text{(new variable)} \\
\text{rchNet}[v][j][i] &= -\text{var\_id} \quad \text{(negation)}
\end{align}

This reduces the number of reachability variables by half!

\subsection{Lazy Capacity Checking}

Instead of encoding capacity constraints upfront (which requires complex cardinality constraints), we check capacity \textbf{after} finding a solution:

\begin{algorithm}[H]
\caption{Check Capacity Violation for CVRP}
\begin{algorithmic}[1]
\Function{CheckCapacity}{$route$, $demands$, $Q$}
    \State $total\_demand \gets \sum_{c \in route} demands[c]$
    \If{$total\_demand > Q$}
        \State \Return (True, GenerateNogood($route$))
    \Else
        \State \Return (False, None)
    \EndIf
\EndFunction
\end{algorithmic}
\end{algorithm}

\subsection{Nogood Clause Generation}

When a route $R = [0, c_1, c_2, \ldots, c_k, 0]$ violates capacity:

\begin{equation}
\text{nogood} = \neg\text{conNet}[v][0][c_1] \vee \neg\text{conNet}[v][c_1][c_2] \vee \cdots \vee \neg\text{conNet}[v][c_k][0]
\end{equation}

This clause says: ``At least one of these edges must NOT be used.''

Adding this clause eliminates the current infeasible solution and allows the solver to find alternatives.

%=============================================================================
\section{Implementation Details}
%=============================================================================

\subsection{Data Structures}

\begin{lstlisting}[caption={Core Data Structures}]
class IncrementalCVRPSolver:
    def __init__(self, n_customers, n_vehicles, distances, demands, capacity):
        self.n = n_customers + 1  # Including depot
        self.m = n_vehicles
        self.distances = distances
        self.demands = demands
        self.capacity = capacity
        
        # Connection network: conNet[v][i][j] = variable ID
        self.conNet = [[[0] * self.n for _ in range(self.n)] 
                       for _ in range(self.m)]
        
        # Reachability network: rchNet[v][i][j] = variable ID (or negation)
        self.rchNet = [[[0] * self.n for _ in range(self.n)] 
                       for _ in range(self.m)]
        
        # WCNF formula
        self.wcnf = WCNF()
\end{lstlisting}

\subsection{Variable Generation}

\begin{lstlisting}[caption={Variable Generation with Optimization}]
def gen_var_for_rch_net(self):
    """Generate reachability variables with negation trick."""
    for v in range(self.m):
        for i in range(1, self.n):
            for j in range(i + 1, self.n):  # Only i < j
                if self.is_valid_edge(v, i, j):
                    var = self.new_var_id()
                    self.rchNet[v][i][j] = var
                    self.rchNet[v][j][i] = -var  # Negation!
\end{lstlisting}

\subsection{Incremental Solving Loop}

\begin{lstlisting}[caption={Main Incremental Solving Loop}]
def solve_incremental(self, timeout=300):
    with RC2(self.wcnf, incr=True) as solver:
        while True:
            model = solver.compute()
            
            if model is None:
                return None  # UNSAT
            
            routes = self.decode_all_routes(model)
            
            # Check capacity for each vehicle
            is_feasible, violation = self.check_solution(routes)
            
            if is_feasible:
                return routes, solver.cost
            else:
                # Generate and add nogood clause
                v, violating_route = violation
                nogood = self.generate_nogood_clause(v, violating_route)
                solver.add_clause(nogood)
\end{lstlisting}

\subsection{K-Nearest Neighbors Integration}

To combine with the CVRP post-improvement approach, we use K-nearest neighbors to restrict the search space:

\begin{lstlisting}[caption={K-Nearest Neighbors for Search Space Reduction}]
def compute_k_neighbors(self, initial_routes=None):
    """Compute valid edges based on K-nearest neighbors."""
    for i in range(self.n):
        # Sort neighbors by distance
        dists = [(j, self.distances[i, j]) 
                 for j in range(self.n) if i != j]
        dists.sort(key=lambda x: x[1])
        
        # Take K nearest
        k_nearest = [j for j, _ in dists[:self.k_neighbors]]
        
        # Mark as valid for all vehicles
        for v in range(self.m):
            for j in k_nearest:
                self.valid_edges[v][i][j] = True
    
    # Ensure initial route edges are valid
    if initial_routes:
        for v, route in initial_routes.items():
            full_route = [0] + route + [0]
            for idx in range(len(full_route) - 1):
                i, j = full_route[idx], full_route[idx + 1]
                self.valid_edges[v][i][j] = True
\end{lstlisting}

%=============================================================================
\section{Complete Algorithm}
%=============================================================================

\begin{algorithm}[H]
\caption{Incremental MaxSAT for CVRP}
\begin{algorithmic}[1]
\Require Instance $(n, m, D, q, Q)$, K-neighbors parameter
\Ensure Optimal routes or UNSAT

\State \textbf{Phase 1: Initialization}
\State Compute K-nearest neighbor matrix
\State Generate connection variables $\text{conNet}[v][i][j]$
\State Generate reachability variables $\text{rchNet}[v][i][j]$

\State \textbf{Phase 2: Generate Base Clauses}
\State Generate soft clauses (objective function)
\State Generate hard clauses (Hamiltonian path constraints)
\State Generate hard clauses (depot in/out constraints)
\State Generate hard clauses (flow conservation)

\State \textbf{Phase 3: Incremental Solving}
\State Initialize RC2 solver with WCNF
\Repeat
    \State $model \gets$ solver.compute()
    \If{$model = $ None}
        \State \Return UNSAT
    \EndIf
    
    \State $routes \gets$ DecodeRoutes($model$)
    \State $feasible, violation \gets$ CheckCapacity($routes$)
    
    \If{NOT $feasible$}
        \State $nogood \gets$ GenerateNogood($violation$)
        \State solver.add\_clause($nogood$)
    \EndIf
\Until{$feasible$}

\State \Return $routes$, solver.cost
\end{algorithmic}
\end{algorithm}

%=============================================================================
\section{Theoretical Analysis}
%=============================================================================

\subsection{Correctness}

\begin{theorem}[Soundness]
If the algorithm returns a solution, it is a feasible CVRP solution.
\end{theorem}

\begin{proof}
The hard clauses enforce:
\begin{enumerate}
    \item Each vehicle leaves and returns to depot exactly once
    \item Each customer is visited exactly once
    \item No subtours (enforced by reachability constraints)
\end{enumerate}
The lazy checking ensures capacity constraints are satisfied before returning.
\end{proof}

\begin{theorem}[Completeness]
If a feasible solution exists, the algorithm will find one.
\end{theorem}

\begin{proof}
The nogood clauses only eliminate infeasible solutions. Since there are finitely many route configurations, the algorithm terminates.
\end{proof}

\subsection{Complexity Analysis}

\begin{itemize}
    \item \textbf{Variables}: $O(m \cdot n^2)$ for connection and reachability
    \item \textbf{Hard clauses}: $O(m \cdot n^3)$ for transitivity
    \item \textbf{Soft clauses}: $O(m \cdot K \cdot n)$ with K-neighbors
    \item \textbf{Iterations}: Worst case exponential, but typically few in practice
\end{itemize}

\subsection{Advantages over MTZ Formulation}

\begin{enumerate}
    \item \textbf{No auxiliary integer variables}: MTZ requires $u_i$ position variables
    \item \textbf{Tighter relaxation}: Reachability constraints provide stronger propagation
    \item \textbf{Incremental solving}: Can add clauses without rebuilding the model
    \item \textbf{Variable reduction}: Negation trick halves reachability variables
\end{enumerate}

%=============================================================================
\section{Comparison with Original CVRP Approach}
%=============================================================================

\begin{table}[H]
\centering
\begin{tabular}{|l|p{5cm}|p{5cm}|}
\hline
\textbf{Feature} & \textbf{Original CVRP (clasp)} & \textbf{Incremental (RTSS-style)} \\
\hline
Solver & clasp (external) & PySAT RC2 (Python) \\
\hline
Subtour Elimination & MTZ with binary encoding & Reachability variables \\
\hline
Capacity Handling & Cardinality constraints & Lazy checking + nogoods \\
\hline
Incremental & No (rebuild for each change) & Yes (add clauses on-the-fly) \\
\hline
Initial Solution & Required (CW + 2-opt) & Optional (warm start) \\
\hline
Variable Count & $O(n^2 \log n)$ for MTZ & $O(n^2)$ for reachability \\
\hline
\end{tabular}
\caption{Comparison of CVRP Solving Approaches}
\end{table}

%=============================================================================
\section{Usage Example}
%=============================================================================

\begin{lstlisting}[caption={Complete Usage Example}]
from incremental_cvrp_solver import IncrementalCVRPSolver
import numpy as np

# Create instance
n_customers = 5
n_vehicles = 2
capacity = 10
demands = np.array([0, 3, 4, 3, 4, 3])  # depot + customers
distances = np.array([...])  # Distance matrix

# Create solver
solver = IncrementalCVRPSolver(
    n_customers=n_customers,
    n_vehicles=n_vehicles,
    distances=distances,
    demands=demands,
    capacity=capacity,
    k_neighbors=4
)

# Solve
routes, cost, time = solver.build_and_solve(timeout=60)

# Output
print(f"Cost: {cost}")
for v, route in routes.items():
    print(f"Vehicle {v}: 0 -> {' -> '.join(map(str, route))} -> 0")
\end{lstlisting}

%=============================================================================
\section{Experimental Results}
%=============================================================================

\subsection{Small Instance Test}

Testing on a 5-customer, 2-vehicle instance:
\begin{itemize}
    \item Customers: 5
    \item Vehicles: 2
    \item Capacity: 10
    \item Demands: [3, 4, 3, 4, 3]
\end{itemize}

\textbf{Results:}
\begin{verbatim}
Cost: 105
Time: 1.33s
Iterations: 21
Capacity violations found: 20

Routes:
  Vehicle 0: 0 -> 5 -> 4 -> 3 -> 0 (Demand: 10/10)
  Vehicle 1: 0 -> 2 -> 1 -> 0 (Demand: 7/10)
\end{verbatim}

The solver correctly found a feasible solution after eliminating 20 capacity-violating solutions through nogood clauses.

%=============================================================================
\section{Conclusion}
%=============================================================================

This document presented a detailed adaptation of the RTSS Incremental MaxSAT approach to CVRP:

\begin{enumerate}
    \item \textbf{Variable Encoding}: Connection variables $l(i,j)$ and reachability variables $r(i,j)$ from RTSS are directly applicable to CVRP.
    
    \item \textbf{Constraint Adaptation}: The main changes are:
    \begin{itemize}
        \item Modified depot constraints (single depot, return required)
        \item Simplified customer model (no pick-up/drop-off pairs)
        \item Removed deadline constraints (capacity only)
    \end{itemize}
    
    \item \textbf{Lazy Checking}: Capacity constraints are checked post-hoc and violations are eliminated via nogood clauses.
    
    \item \textbf{Incremental Solving}: The RC2 solver supports adding clauses incrementally, enabling efficient exploration of the solution space.
    
    \item \textbf{K-Neighbors Integration}: The search space can be restricted using K-nearest neighbors, combining the benefits of both approaches.
\end{enumerate}

The implementation is provided in \texttt{incremental\_cvrp\_solver.py}.

\end{document}
